\chapter{Building the \cotwo{} dataset}
\label{ch:dataset}

\section{Why the \cotwo{} is computed separately}
The forecast model needs the \cotwo{} for 30~minutes per setting, as in \citet{bian2025}: the \cotwo{} is
still rising after 10~minutes when the ventilation is weak, and 30~minutes give 44 training samples of
9~minutes each (6~minutes of history and 3~minutes of forecast). A 30-minute OpenFOAM run would take about
two days per setting on one core. I therefore split the work into two parts:
\begin{enumerate}
  \item OpenFOAM computes the airflow for 600~s and saves the velocity and the turbulent viscosity.
  \item My own solver computes the \cotwo{} for the full 30~minutes on the GPU from this airflow, and saves it
        every 30~s.
\end{enumerate}
When the windows open, the airflow needs about 2.5~minutes to develop. After that it fluctuates around the
same pattern, because the window setting does not change during a run. The \cotwo{} solver therefore uses the
real OpenFOAM flow for the first 150~s. The flow saved between 150 and 600~s is averaged into one mean flow
field, and from 160~s on the solver keeps this mean flow fixed until 30~minutes (the change from the last
real flow to the mean flow takes 10~s). Only the airflow is kept fixed; the \cotwo{} is computed at every
time step.

\section{The \cotwo{} solver}
The solver works on the same 8~cm grid as OpenFOAM, on the GPU (PyTorch). Before use, every flow field is
made divergence-free: the velocities on the cell faces are corrected with a pressure-like field from a
Poisson equation (conjugate gradients), so that exactly as much air leaves through the open door as comes in
through the windows. The advection is computed in flux form with a minmod-limited upwind scheme, so the
\cotwo{} never becomes negative and the total amount of \cotwo{} changes only through the source and the
outflow. The time integration is a three-stage strong-stability-preserving Runge--Kutta scheme with a time
step limited by the Courant number (0.5) and by the diffusion. The diffusivity is $D+\nu_t/\mathrm{Sc}_t$
(\eref{eq:co2}). Fresh air enters with \cotwo{}~$=0$ (outdoor level), and \cotwo{} leaves through the door
with the outflowing air. \todo{short validation of the solver itself: closed room / mass balance tests}

\section{Checks against OpenFOAM}
OpenFOAM computes the \cotwo{} itself on its own flow. I use this as the reference and compare it with my
\cotwo{} at the last OpenFOAM time, in two ways: as the relative $\ell_2$ difference of the \cotwo{} above
the outdoor level in the whole room (strict), and in the metric of \citet{bian2025}, which includes the
400~ppm of the outdoor air and is evaluated on the planes at 1.1 and 1.6~m.

In the first setup the checks showed two weaknesses: the total amount of \cotwo{} was not kept exactly, and
the mean flow mixed the \cotwo{} too little, because the averaging removes the fluctuations. After making the
solver mass-conserving and using a turbulent Schmidt number of $\mathrm{Sc}_t=0.3$ instead of 0.7
\citep{tominaga2007}, the difference in the paper's metric fell from 6--9\,\% to 1--4\,\%. In one setting
that was simulated in OpenFOAM for the full 30~minutes, my \cotwo{} stayed within about 3\,\% of OpenFOAM's
over the whole time.

\section{Reproducibility and simulated time}
To check that the data can be reproduced, I simulated setting N10 a second time
and built its dataset again. The \cotwo{} of both runs differs by 1.1\,\% over all points and times, and by
at most 0.15\,\% in the paper's metric. Using only the first 450~s of the OpenFOAM run instead of 600~s
changes the \cotwo{} by 3.1\,\% (paper's metric at most 0.48\,\%) and would make the simulations about
1.33 times faster.

\section{The save interval of the flow}
Comparing all 28 settings with OpenFOAM's \cotwo{} showed large differences in some settings with high
window speeds, already when my solver used the real OpenFOAM flow without any averaging (for example
114\,\% in N14 and 131\,\% in N28 after 600~s). The settings with these differences all have W1 and W3 open
at high speed and a strongly fluctuating flow.

The reason was the save interval. OpenFOAM saved the flow every 10~s, and my solver interpolates linearly
between the saves. At 2~m/s the air moves about 20~m in 10~s, more than the length of the room, so the saves
miss most of what the air does. To test this, I simulated N28 for 150~s with a save every second and computed
the \cotwo{} from every 1st, 2nd, 5th and 10th save (\tref{tab:saves}).

\begin{table}[h]
  \centering
  \caption{N28: my \cotwo{} against OpenFOAM's \cotwo{} for different save intervals of the flow (relative
  $\ell_2$ difference in the whole room / total \cotwo{} ours : OpenFOAM).}
  \label{tab:saves}
  \footnotesize
  \begin{tabular}{lrrrrr}
    \toprule
    Saved every & 30~s & 60~s & 90~s & 120~s & 150~s\\
    \midrule
    1~s  & 3.1\,\% / 1.001 & 2.1\,\% / 0.998 & 2.2\,\% / 0.998 & 2.0\,\% / 1.000 & 2.0\,\% / 1.001\\
    2~s  & 4.3\,\% / 1.004 & 4.4\,\% / 0.999 & 4.0\,\% / 1.002 & 4.0\,\% / 1.008 & 4.7\,\% / 1.012\\
    5~s  & 12.5\,\% / 1.021 & 16.0\,\% / 1.002 & 13.1\,\% / 1.027 & 17.7\,\% / 1.077 & 19.4\,\% / 1.103\\
    10~s & 30.6\,\% / 1.032 & 29.7\,\% / 0.997 & 27.2\,\% / 1.060 & 40.0\,\% / 1.156 & 45.4\,\% / 1.232\\
    \bottomrule
  \end{tabular}
\end{table}

With saves every second my solver matches OpenFOAM within about 2\,\%. N14 gives the same picture: 1.9--3.7\,\%
with 1~s saves and 24--45\,\% with 10~s saves. For the calm setting N01 (one window at 0.5~m/s) the save
interval hardly matters (25.8\,\% with 10~s and 25.6\,\% with 1~s saves at 150~s, both with $\mathrm{Sc}_t=0.3$).
So my solver itself is accurate; only its input was too coarse in fast settings.

\section{The flow after 150~s}
With a run of 450~s and saves every second for N28 and N01, I also tested how the \cotwo{} after 150~s can be
made more accurate (\tref{tab:after150}). The real flow saved every second gives about 2\,\% at all times,
which is the best that is possible with this solver.

\begin{table}[h]
  \centering
  \caption{Methods for the \cotwo{} after 150~s, against OpenFOAM's \cotwo{} at 450~s (relative $\ell_2$
  difference in the whole room; in brackets the paper's metric at 1.1~m).}
  \label{tab:after150}
  \small
  \begin{tabular}{>{\raggedright}p{8.2cm}rr}
    \toprule
    Method & N28 (fast) & N01 (calm)\\
    \midrule
    10~s saves to 150~s, then mean flow, $\mathrm{Sc}_t=0.3$ (first dataset) & 106.9\,\% (3.4\,\%) & 32.8\,\% (10.1\,\%)\\
    1~s saves to 150~s, then mean flow, $\mathrm{Sc}_t=0.3$ & 100.7\,\% (3.2\,\%) & 33.1\,\% (10.2\,\%)\\
    1~s saves to 150~s, then mean flow, $\mathrm{Sc}_t=0.7$ & 143.6\,\% (5.6\,\%) & 35.2\,\% (11.9\,\%)\\
    real flow to 300~s, then the 150--300~s flow in a loop, $\mathrm{Sc}_t=0.7$ & 42.2\,\% (1.5\,\%) & 76.0\,\% (31.6\,\%)\\
    real flow to 450~s (upper limit), $\mathrm{Sc}_t=0.7$ & 1.9\,\% (0.1\,\%) & 1.4\,\% (0.3\,\%)\\
    \bottomrule
  \end{tabular}
\end{table}

Neither alternative works for both settings. In the fast setting the fixed mean flow keeps too much
\cotwo{} in the room (about 1.8 times OpenFOAM's amount in N28 at 450~s): the mean flow has calm zones that trap the
\cotwo{}, while the real, fluctuating flow flushes them. Repeating the real flow in a loop helps there, but in
the calm setting it is much worse, probably because the slow flow keeps changing and the first part is not
typical for the later part. A turbulent Schmidt number of 0.7 is right while the real flow is used
(1.9\,\% at 150~s instead of 25.6\,\% in N01), but with the mean flow 0.3 is better. I therefore keep the mean
flow after 150~s and treat its error as a known limitation. The only way to remove it would be to simulate
the full 30~minutes in OpenFOAM.

\section{The corrected dataset}
Based on these tests I rebuilt the dataset (data~v3) in the same way for all 28 settings:
\begin{itemize}
  \item OpenFOAM is run again for the first 150~s with a save every second;
  \item for 0--150~s my solver uses this flow with $\mathrm{Sc}_t=0.7$, as OpenFOAM;
  \item after 150~s it uses the mean flow of the first dataset with $\mathrm{Sc}_t=0.3$.
\end{itemize}
At 150~s the new dataset matches OpenFOAM's \cotwo{} within 1.0--2.2\,\% in the settings I have checked so far (paper's
metric at most 0.1\,\%), compared with up to 45\,\% before. \todo{range for all 28 settings from the stored checks}

\tref{tab:v3v2} shows how much the data changed. The differences are in the first minutes, when the fresh air
spreads into the room. After about 10~minutes both datasets become the same, because they use the same mean
flow, and at 30~minutes they are identical. Since the difference after 20~minutes is at most 0.15\,\%, more
than ten times smaller than the accuracy of the solver, the last 10~minutes of 20 of the 28 settings were taken from
the first dataset to save computing time.

\begin{table}[h]
  \centering
  \caption{Difference between the corrected and the first dataset (relative $\ell_2$ difference of the
  \cotwo{} in the whole room), for the first 8 settings that were finished.}
  \label{tab:v3v2}
  \begin{tabular}{lrrrrr}
    \toprule
    Setting & 1~min & 2.5~min & 5~min & 10~min & 30~min\\
    \midrule
    N01 & 24.8\,\% & 29.3\,\% & 1.4\,\% & 0.2\,\% & 0.0\,\%\\
    N07 & 32.1\,\% & 26.3\,\% & 2.1\,\% & 0.6\,\% & 0.1\,\%\\
    N09 & 34.5\,\% & 18.5\,\% & 1.5\,\% & 0.2\,\% & 0.0\,\%\\
    N10 & 36.8\,\% & 31.1\,\% & 1.8\,\% & 0.4\,\% & 0.1\,\%\\
    N14 & 38.1\,\% & 35.9\,\% & 5.1\,\% & 0.7\,\% & 0.0\,\%\\
    N20 & 29.2\,\% & 23.8\,\% & 0.9\,\% & 0.2\,\% & 0.0\,\%\\
    N27 & 43.8\,\% & 27.2\,\% & 1.9\,\% & 0.4\,\% & 0.0\,\%\\
    N28 & 36.3\,\% & 36.1\,\% & 7.8\,\% & 1.6\,\% & 0.0\,\%\\
    \bottomrule
  \end{tabular}
\end{table}

\section{Paper-like transition data}
In \citet{bian2025} every run starts from the steady state of another setting. I do the same without new
OpenFOAM runs: for each setting I start from the \cotwo{} of a randomly chosen other setting after 30~minutes,
switch off the people, and let the new flow carry this \cotwo{} away. Because \eref{eq:co2} is linear, the
\cotwo{} of the new setting's people can then be added with any number of people, and the numbers before and
after the change are drawn between 10 and 80 during training, as in the paper.

For the corrected dataset the transition runs had to be computed again for the full 30~minutes. Here the
\cotwo{} only decays, so the difference to the old transition runs does not become small over time; for the
runs computed so far it is 2.6--13.6\,\% at 30~minutes. \todo{final range for all 28 settings}
